\documentclass[../../main.tex]{subfiles}
\begin{document}

\begin{center}
  {\fontsize{48pt}{58pt}\selectfont \textbf{1967}\normalsize\textbf{年度}}
\end{center}

\begin{table}[htb]
  \centering
  \caption{出題テーマと難易度}
  \begin{tabular}{|c|p{6.2cm}|c|c|c|}
    \hline
    \textbf{問題番号} & \textbf{テーマ} & \textbf{思考量} & \textbf{計算量} & \textbf{総合} \\
    \hline
    第1問 & 相加相乗平均による不等式 & 易 & 易 & 易 \\
    \hline
    第2問 & 飛行機の方向と仰角から求める高度 & 普 & 普 & 普 \\
    \hline
    第3問 & 動く正方形と固定正方形の共通面積 & 普 & 普 & 普 \\
    \hline
    第4問 & 楕円で囲まれる領域の面積 & 普 & 普 & 普 \\
    \hline
    第5問 & 新課程：放物線と対数曲線／旧課程：三角関数の接点 & 易 & 易 & 易 \\
    \hline
    第6問 & 新課程：番号が一致するカード／旧課程：半円から定まる点列 & 易 & 易 & 易 \\
    \hline
  \end{tabular}
\end{table}

\begin{table}[htb]
  \centering
  \caption{解答}
  \renewcommand{\arraystretch}{1.4}
  \begin{tabular}{|c|c|p{9cm}|}
    \hline
    \textbf{問題番号} & \textbf{小問} & \textbf{答え} \\
    \hline
    第1問 & -- & $ax^{n+1}+\dfrac{1}{\sqrt[n]{a}}>x$（証明問題） \\
    \hline
    第2問 & -- & 高度は$1000\,\mathrm{m}$ \\
    \hline
    第3問 & -- & 最大面積は$\dfrac{3\sqrt{3}}{4}$ \\
    \hline
    第4問 & -- & 面積は$\dfrac{2\sqrt{3}\pi}{3}$（略図は解答参照） \\
    \hline
    第5問 & 新課程 & $a=\dfrac{1}{2e}$．面積は$\dfrac{2\sqrt{e}}{3}-1$ \\
    \cline{2-3}
    & 旧課程 & $a=-3,1,\dfrac{3}{2}$（グラフは解答参照） \\
    \hline
    第6問 & 新課程 & 確率は$\dfrac{1}{320}$ \\
    \cline{2-3}
    & 旧課程 & $A_0A=4$ \\
    \hline
  \end{tabular}
\end{table}

\paragraph{試験形態}
6問．第5・6問は新旧課程別．試験時間は要確認．

\end{document}
